\documentclass[11pt,a4paper]{article}
\usepackage[utf8]{inputenc}
\usepackage[T1]{fontenc}
\usepackage[french]{babel}
\usepackage{amsmath, amssymb, amsthm}
\usepackage{geometry}
\usepackage{xcolor}
\geometry{hmargin=2cm, vmargin=2cm}

% Macro pour le formatage des raisonnements
\newcommand{\raisonnement}[1]{
    \vspace{0.2cm}
    \noindent\textcolor{blue!70!black}{
        \textbf{\textit{Raisonnement : }} \textit{#1}
    }
    \vspace{0.2cm}
}

\begin{document}

\begin{center}
    \textbf{\Large Corrigé de l'Examen MA333 - 2023-2024} \\
    \vspace{0.2cm}
    \textbf{GRENOBLE INP Esisar UGA, VALENCE}
\end{center}

\vspace{0.5cm}

% ==========================================
% EXERCICE 1
% ==========================================
\section*{Exercice 1 (4 points)}

Soit $X \sim \mathcal{N}(75, 4^2)$. On pose $Z = \frac{X-75}{4}$ la variable centrée réduite associée, avec $Z \sim \mathcal{N}(0,1)$. On note $\Phi$ la fonction de répartition de la loi normale centrée réduite.

\textbf{1. Calculer $\mathbb{P}(X>80)$ et $\mathbb{P}(65<X<85)$.}
\raisonnement{On passe à la variable centrée réduite $Z$ pour pouvoir lire les probabilités dans la table de la loi normale.}
\begin{align*}
\mathbb{P}(X > 80) &= \mathbb{P}\left(Z > \frac{80-75}{4}\right) = \mathbb{P}(Z > 1.25) \\
&= 1 - \Phi(1.25) \approx 1 - 0.8944 = \mathbf{0.1056} \\
\mathbb{P}(65 < X < 85) &= \mathbb{P}\left(\frac{65-75}{4} < Z < \frac{85-75}{4}\right) = \mathbb{P}(-2.5 < Z < 2.5) \\
&= 2\Phi(2.5) - 1 \approx 2(0.9938) - 1 = \mathbf{0.9876}
\end{align*}

\textbf{2. Déterminer la loi de $S$.}
\raisonnement{La somme de $n$ variables aléatoires indépendantes suivant des lois normales suit elle-même une loi normale. L'espérance de la somme est la somme des espérances, et par indépendance, la variance de la somme est la somme des variances.}
On a $S = \sum_{i=1}^n X_i$. 
$$ \mathbb{E}(S) = \sum_{i=1}^n \mathbb{E}(X_i) = n \times 75 = 75n $$
$$ \mathbb{V}(S) = \sum_{i=1}^n \mathbb{V}(X_i) = n \times 4^2 = 16n $$
Donc, \textbf{$S$ suit une loi normale $\mathcal{N}(75n, 16n)$} (d'écart-type $4\sqrt{n}$).

\textbf{3. Nombre maximum de personnes pour un risque de surcharge $\le 10^{-3}$.}
\raisonnement{On cherche le plus grand entier $n$ tel que $\mathbb{P}(S > 500) \le 10^{-3}$. On utilise à nouveau le centrage et la réduction, appliqués cette fois à $S$.}
\begin{align*}
\mathbb{P}(S > 500) \le 0.001 &\iff \mathbb{P}\left(\frac{S - 75n}{4\sqrt{n}} > \frac{500 - 75n}{4\sqrt{n}}\right) \le 0.001 \\
&\iff 1 - \Phi\left(\frac{500 - 75n}{4\sqrt{n}}\right) \le 0.001 \\
&\iff \Phi\left(\frac{500 - 75n}{4\sqrt{n}}\right) \ge 0.999
\end{align*}
D'après la table de la loi normale, $\Phi(3.09) \approx 0.999$. Il faut donc :
$$ \frac{500 - 75n}{4\sqrt{n}} \ge 3.09 $$
Testons les valeurs de $n$ :
\begin{itemize}
    \item Si $n=6$, $\mu = 450$ et $\sigma = 4\sqrt{6} \approx 9.798$. On a $z = \frac{50}{9.798} \approx 5.10 \ge 3.09$. Le risque est largement inférieur à $10^{-3}$.
    \item Si $n=7$, $\mu = 525$. L'espérance dépasse la charge limite, le risque de surcharge sera supérieur à $50\%$.
\end{itemize}
Le nombre maximum de personnes autorisées est donc de \textbf{6 personnes}.

\vspace{0.5cm}

% ==========================================
% EXERCICE 2
% ==========================================
\section*{Exercice 2 (4 points)}

Soit $f(x,y) = 4y(1-x)\mathbb{I}_{[0;1]}(x)\mathbb{I}_{[0;1]}(y)$.

\textbf{1. Montrer que $f$ est une densité de probabilité.}
\raisonnement{Il faut vérifier deux conditions : $f$ doit être positive ou nulle sur $\mathbb{R}^2$, et l'intégrale double de $f$ sur $\mathbb{R}^2$ doit être égale à 1.}
\begin{itemize}
    \item Pour tout $(x,y) \in [0,1]^2$, $y \ge 0$ et $1-x \ge 0$, donc $f(x,y) \ge 0$. Ailleurs, $f(x,y) = 0$.
    \item Calculons l'intégrale sur $\mathbb{R}^2$ :
    $$ \iint_{\mathbb{R}^2} f(x,y) dx dy = \int_0^1 \int_0^1 4y(1-x) dx dy = 4 \left( \int_0^1 (1-x) dx \right) \left( \int_0^1 y dy \right) $$
    $$ = 4 \left[ x - \frac{x^2}{2} \right]_0^1 \left[ \frac{y^2}{2} \right]_0^1 = 4 \times \frac{1}{2} \times \frac{1}{2} = 1 $$
\end{itemize}
$f$ est bien une densité de probabilité.

\textbf{2. Calculer $\mathbb{P}(X < Y)$.}
\raisonnement{Il s'agit d'intégrer la densité sur le domaine défini par $0 \le x < y \le 1$.}
\begin{align*}
\mathbb{P}(X < Y) &= \int_0^1 \left( \int_0^y 4y(1-x) dx \right) dy \\
&= \int_0^1 4y \left[ x - \frac{x^2}{2} \right]_0^y dy \\
&= \int_0^1 4y \left( y - \frac{y^2}{2} \right) dy = \int_0^1 (4y^2 - 2y^3) dy \\
&= \left[ \frac{4y^3}{3} - \frac{2y^4}{4} \right]_0^1 = \frac{4}{3} - \frac{1}{2} = \mathbf{\frac{5}{6}}
\end{align*}

\textbf{3. Déterminer les lois de X et de Y.}
\raisonnement{Pour obtenir les densités marginales, on intègre la densité conjointe par rapport à l'autre variable sur son domaine de définition.}
\begin{align*}
f_X(x) &= \int_{-\infty}^{+\infty} f(x,y) dy = \left( \int_0^1 4y(1-x) dy \right) \mathbb{I}_{[0;1]}(x) = 4(1-x) \left[ \frac{y^2}{2} \right]_0^1 \mathbb{I}_{[0;1]}(x) = \mathbf{2(1-x) \mathbb{I}_{[0;1]}(x)} \\
f_Y(y) &= \int_{-\infty}^{+\infty} f(x,y) dx = \left( \int_0^1 4y(1-x) dx \right) \mathbb{I}_{[0;1]}(y) = 4y \left[ x - \frac{x^2}{2} \right]_0^1 \mathbb{I}_{[0;1]}(y) = \mathbf{2y \mathbb{I}_{[0;1]}(y)}
\end{align*}

\textbf{4. Les variables aléatoires X et Y sont-elles indépendantes ?}
\raisonnement{Deux variables à densité sont indépendantes si et seulement si leur densité conjointe est le produit de leurs densités marginales presque partout.}
On vérifie si $f(x,y) = f_X(x) f_Y(y)$ :
$$ f_X(x) \times f_Y(y) = 2(1-x)\mathbb{I}_{[0;1]}(x) \times 2y\mathbb{I}_{[0;1]}(y) = 4y(1-x)\mathbb{I}_{[0;1]}(x)\mathbb{I}_{[0;1]}(y) = f(x,y) $$
Oui, \textbf{$X$ et $Y$ sont indépendantes}.

\vspace{0.5cm}

% ==========================================
% EXERCICE 3
% ==========================================
\section*{Exercice 3 (6 points)}

Pour la pièce A, on définit le "succès" comme l'obtention de "face", de probabilité $1-a$. Pour la pièce B, le succès est "face" avec probabilité $1-b$.

\textbf{1. Quelles sont les lois de probabilités de X et de Y ? Donner $\mathbb{E}(X)$ et $\mathbb{E}(Y)$.}
\raisonnement{X et Y représentent le nombre d'essais nécessaires pour obtenir le premier succès lors d'une succession d'épreuves de Bernoulli indépendantes. Elles suivent donc des lois géométriques.}
\textbf{$X \sim \mathcal{G}(1-a)$} et \textbf{$Y \sim \mathcal{G}(1-b)$}.
Pour tout $k \in \mathbb{N}^*$, $\mathbb{P}(X=k) = a^{k-1}(1-a)$ et $\mathbb{P}(Y=k) = b^{k-1}(1-b)$. \\
Les espérances sont : $\mathbb{E}(X) = \mathbf{\frac{1}{1-a}}$ et $\mathbb{E}(Y) = \mathbf{\frac{1}{1-b}}$.

\textbf{2. Calculer la probabilité de l'évènement $(X=Y)$.}
\raisonnement{Les lancers des deux pièces étant indépendants, les variables $X$ et $Y$ le sont aussi. On somme sur tous les cas possibles $k$.}
\begin{align*}
\mathbb{P}(X=Y) &= \sum_{k=1}^{+\infty} \mathbb{P}(X=k \cap Y=k) = \sum_{k=1}^{+\infty} \mathbb{P}(X=k)\mathbb{P}(Y=k) \\
&= \sum_{k=1}^{+\infty} a^{k-1}(1-a) b^{k-1}(1-b) = (1-a)(1-b) \sum_{k=1}^{+\infty} (ab)^{k-1} \\
&= (1-a)(1-b) \sum_{j=0}^{+\infty} (ab)^{j} = \mathbf{\frac{(1-a)(1-b)}{1-ab}}
\end{align*}

\textbf{3. Trouver, pour $k \in \mathbb{N}$, la valeur de $\mathbb{P}(X>k)$. En déduire $\mathbb{P}(X>Y)$ et $\mathbb{P}(X \ge Y)$.}
\raisonnement{Dire que l'on a besoin de plus de $k$ lancers pour obtenir le premier "face" revient à dire que les $k$ premiers lancers ont tous donné "pile".}
$$ \mathbb{P}(X>k) = \mathbf{a^k} $$
Pour $\mathbb{P}(X>Y)$, on utilise la formule des probabilités totales en fixant la valeur de $Y$ :
\begin{align*}
\mathbb{P}(X>Y) &= \sum_{k=1}^{+\infty} \mathbb{P}(Y=k) \mathbb{P}(X>k) = \sum_{k=1}^{+\infty} b^{k-1}(1-b) a^k = a(1-b) \sum_{k=1}^{+\infty} (ab)^{k-1} = \mathbf{\frac{a(1-b)}{1-ab}}
\end{align*}
Pour $\mathbb{P}(X \ge Y)$ :
$$ \mathbb{P}(X \ge Y) = \mathbb{P}(X>Y) + \mathbb{P}(X=Y) = \frac{a(1-b) + (1-a)(1-b)}{1-ab} = \frac{(1-b)(a + 1 - a)}{1-ab} = \mathbf{\frac{1-b}{1-ab}} $$

\textbf{4. Calculer la probabilité $\mathbb{P}(M \ge k)$. En déduire la loi de M.}
\raisonnement{Le minimum de deux variables est supérieur ou égal à $k$ si et seulement si les deux variables sont simultanément supérieures ou égales à $k$.}
Pour $k \in \mathbb{N}^*$, $\mathbb{P}(M \ge k) = \mathbb{P}(X \ge k \cap Y \ge k) = \mathbb{P}(X \ge k)\mathbb{P}(Y \ge k)$ par indépendance.
Or $\mathbb{P}(X \ge k) = \mathbb{P}(X > k-1) = a^{k-1}$.
Donc :
$$ \mathbb{P}(M \ge k) = a^{k-1}b^{k-1} = \mathbf{(ab)^{k-1}} $$
On en déduit la loi ponctuelle de $M$ :
$$ \mathbb{P}(M=k) = \mathbb{P}(M \ge k) - \mathbb{P}(M \ge k+1) = (ab)^{k-1} - (ab)^k = \mathbf{(ab)^{k-1}(1-ab)} $$
On reconnaît la formule d'une \textbf{loi géométrique de paramètre $1-ab$} : $M \sim \mathcal{G}(1-ab)$.

\vspace{0.5cm}

% ==========================================
% EXERCICE 4
% ==========================================
\section*{Exercice 4 (16 points - note: barème original)}

\textbf{1(a). Calculer la probabilité de l'évènement "Robert tombe sur la piste qu'il a choisie".}
\raisonnement{On note $B, R, N$ les évènements de choix de piste et $T$ l'évènement "tomber". On applique la formule des probabilités totales formant un système complet d'évènements.}
Données : $\mathbb{P}(B)=1/4$, $\mathbb{P}(R)=1/2$, $\mathbb{P}(N)=1/4$. Et $\mathbb{P}(T|B)=1/10$, $\mathbb{P}(T|R)=1/6$, $\mathbb{P}(T|N)=2/5$.
\begin{align*}
\mathbb{P}(T) &= \mathbb{P}(T|B)\mathbb{P}(B) + \mathbb{P}(T|R)\mathbb{P}(R) + \mathbb{P}(T|N)\mathbb{P}(N) \\
&= \left(\frac{1}{10} \times \frac{1}{4}\right) + \left(\frac{1}{6} \times \frac{1}{2}\right) + \left(\frac{2}{5} \times \frac{1}{4}\right) = \frac{1}{40} + \frac{1}{12} + \frac{1}{10} \\
&= \frac{3 + 10 + 12}{120} = \frac{25}{120} = \mathbf{\frac{5}{24}}
\end{align*}

\textbf{1(b). Sachant cela, quelle est la probabilité que Robert ait choisi la piste noire ?}
\raisonnement{C'est une inversion de probabilité classique, on utilise la formule de Bayes.}
$$ \mathbb{P}(N|T) = \frac{\mathbb{P}(T|N)\mathbb{P}(N)}{\mathbb{P}(T)} = \frac{\frac{2}{5} \times \frac{1}{4}}{\frac{5}{24}} = \frac{\frac{1}{10}}{\frac{5}{24}} = \frac{1}{10} \times \frac{24}{5} = \frac{24}{50} = \mathbf{\frac{12}{25}} = 0.48 $$

\textbf{2(a). Soit $Z=X+Y$, déterminer la loi de Z.}
\raisonnement{La loi de la somme de deux variables aléatoires indépendantes continues s'obtient par le produit de convolution de leurs densités : $f_Z(z) = \int_{-\infty}^{+\infty} f_X(z-y)f_Y(y) dy$.}
On a $X \sim \mathcal{E}(1/5)$ et $Y \sim \mathcal{U}[10; 15]$.
\begin{align*}
f_Z(z) &= \int_{10}^{15} \frac{1}{5}e^{-\frac{1}{5}(z-y)} \mathbb{I}_{\{z-y \ge 0\}} \times \frac{1}{5} dy = \frac{1}{25}e^{-\frac{z}{5}} \int_{10}^{15} e^{\frac{y}{5}} \mathbb{I}_{\{y \le z\}} dy
\end{align*}
Trois cas se présentent :
\begin{itemize}
    \item Si $z < 10$ : l'intervalle d'intégration est vide, $\mathbf{f_Z(z) = 0}$.
    \item Si $10 \le z \le 15$ : on intègre $y$ de 10 à $z$.
    $$ f_Z(z) = \frac{1}{25}e^{-\frac{z}{5}} \left[ 5e^{\frac{y}{5}} \right]_{10}^z = \frac{1}{5}e^{-\frac{z}{5}} \left( e^{\frac{z}{5}} - e^2 \right) = \mathbf{\frac{1}{5} \left( 1 - e^{-\frac{z-10}{5}} \right)} $$
    \item Si $z > 15$ : on intègre $y$ de 10 à 15.
    $$ f_Z(z) = \frac{1}{25}e^{-\frac{z}{5}} \left[ 5e^{\frac{y}{5}} \right]_{10}^{15} = \mathbf{\frac{1}{5} \left( e^3 - e^2 \right) e^{-\frac{z}{5}}} $$
\end{itemize}

\textbf{2(b). Déterminer la probabilité que Robert ne soit pas en retard.}
\raisonnement{Entre 10h45 et 11h, Robert dispose de 15 minutes. Il ne sera pas en retard si $Z \le 15$.}
\begin{align*}
\mathbb{P}(Z \le 15) &= \int_{10}^{15} f_Z(z) dz = \int_{10}^{15} \frac{1}{5} \left( 1 - e^{-\frac{z-10}{5}} \right) dz \\
&= \left[ \frac{z}{5} + e^{-\frac{z-10}{5}} \right]_{10}^{15} = (3 + e^{-1}) - (2 + 1) = \mathbf{e^{-1}} \approx 0.368
\end{align*}

\textbf{2(c). Probabilité que Stéphanie arrive avant Robert.}
\raisonnement{On cherche $\mathbb{P}(T_2 < T_1)$. L'aire totale de la zone des temps possibles est l'aire du rectangle $[55;85] \times [65;75]$, soit $30 \times 10 = 300$. On calcule géométriquement l'aire du sous-domaine où $t_2 < t_1$.}
Le domaine où $t_2 < t_1$ dans le rectangle est composé d'un triangle pour $t_1 \in [65;75]$ (aire = $10 \times 10 / 2 = 50$) et d'un rectangle complet pour $t_1 \in [75;85]$ (aire = $10 \times 10 = 100$).
Aire totale favorable = 150.
$$ \mathbb{P}(T_2 < T_1) = \frac{150}{300} = \mathbf{\frac{1}{2}} $$

\textbf{3(a). Quelle est la loi exacte de $S_n$ ?}
\raisonnement{On compte le nombre de succès (skis Castafiore de probabilité $p$) parmi $n$ skieurs indépendants.}
\textbf{$S_n$ suit une loi binomiale $\mathcal{B}(n, p)$}.

\textbf{3(b). Déterminer $n$ pour une précision à $10^{-2}$ près au risque $0.05$.}
\raisonnement{La longueur de la demi-amplitude de l'intervalle de confiance asymptotique est $1.96 \sqrt{\frac{p(1-p)}{n}}$. Pour garantir cette marge d'erreur sans connaître $p$, on se place dans le cas le plus défavorable : $p=0.5$.}
$$ 1.96 \sqrt{\frac{0.5 \times 0.5}{n}} \le 0.01 \iff \frac{1.96}{2\sqrt{n}} \le 0.01 \iff \sqrt{n} \ge \frac{0.98}{0.01} = 98 $$
En élevant au carré : $n \ge 98^2 = 9604$.
Il faut donc au minimum \textbf{9604 skieurs}.

\textbf{3(c). Intervalle de confiance de $p$ au risque $5\%$.}
\raisonnement{L'estimation ponctuelle est $\hat{p} = \frac{2500}{8500}$. On applique la formule de l'intervalle de confiance usuel pour les proportions.}
$\hat{p} = \frac{25}{85} = \frac{5}{17} \approx 0.2941$.
La marge d'erreur $E$ est $1.96 \sqrt{\frac{\hat{p}(1-\hat{p})}{n}} = 1.96 \sqrt{\frac{\frac{5}{17} \times \frac{12}{17}}{8500}} \approx 1.96 \times 0.00494 \approx 0.0097$.
L'intervalle de confiance est donc :
$$ IC_{95\%} = [\hat{p} - E ; \hat{p} + E] = [0.2941 - 0.0097 ; 0.2941 + 0.0097] = \mathbf{[0.284 ; 0.304]} $$

\end{document}